% GitHub Repo and Documentation: https://github.com/celiobjunior/resume-template
% Copyright © 2025 Celio B Junior. All rights reserved.
% 
% Licensed under the Apache License, Version 2.0 (the "License");
% you may not use this file except in compliance with the License.
% You may obtain a copy of the License at
%
%     http://www.apache.org/licenses/LICENSE-2.0
%
% This template follows best practices from README.md

% Início do documento LaTeX: define tipo e formato do documento
% a4paper = tamanho A4, 10pt = tamanho base da fonte
\documentclass[a4paper,10pt]{article}

% --- PACOTES: BASE / IDIOMA ---
\usepackage[utf8]{inputenc}
\usepackage[T1]{fontenc}
\usepackage[portuguese]{babel}

% --- PACOTES: LAYOUT / TIPOGRAFIA ---
\usepackage{geometry}
\usepackage{parskip}
\usepackage{microtype}
\usepackage{enumitem}
\usepackage{titlesec}
\usepackage{array}
\usepackage{tabularx}

% --- PACOTES: LINKS / BOOKMARKS ---
\usepackage{hyperref}
\usepackage{bookmark}
\usepackage{xurl}

% --- CONFIGURAÇÃO: MARGENS / ESPAÇAMENTO ---
\geometry{top=1.5cm, bottom=1.5cm, left=1.5cm, right=1.5cm} % Margens da página
\setcounter{secnumdepth}{0}                                 % Remove numeração de seção
\setlist[itemize]{
    leftmargin=0.75em, % Recuo da lista (mais perto da margem da página)
    itemsep=0.2em,     % Espaço entre itens
    topsep=0.25em,     % Espaço antes/depois da lista
    parsep=0em,        % Espaço entre parágrafos no mesmo item
    partopsep=0em      % Espaço extra ao iniciar a lista
}

% Remove números de página e cabeçalhos para visual limpo do currículo
\pagestyle{empty}

% --- CONFIGURAÇÃO: METADADOS / LINKS ---
\hypersetup{
    pdftitle={CV Thiago Buarque - Desenvolvedor Backend .NET}, % Título do PDF
    pdfauthor={Thiago Buarque},                                % Autor do PDF
    colorlinks=true,          % Usa cor nos links (sem caixa)
    linkcolor=black,          % Cor de links internos
    urlcolor=black,           % Cor de URLs
    citecolor=black,          % Cor de citações
    bookmarksdepth=2          % Barra lateral: seção e subseção
}

% --- CONFIGURAÇÃO: TÍTULOS ---
\titleformat{\section}
{\Large\bfseries}
{}
{0em}
{}
[\titlerule\vspace{0.5ex}]
\titleformat{\subsection}
{\normalsize\bfseries}
{}
{0em}
{}
\titlespacing*{\subsection}{0pt}{0.35em}{0.15em}

% Macro para entradas do currículo com bookmark no PDF
\newcounter{cventry}
\newcommand{\cventry}[4]{%
    \refstepcounter{cventry}%
    \phantomsection%
    \pdfbookmark[2]{#1}{cventry-\thecventry}%
    \noindent\begin{tabularx}{\textwidth}{@{}>{\raggedright\arraybackslash}X >{\raggedleft\arraybackslash}X@{}}
    \textbf{#1} & #2 \\
    \textit{#3} & \textit{#4} \\
    \end{tabularx}
}

% Macro para projetos: nome (com link) / stack / descrição
\newcommand{\projectentry}[3]{%
    \phantomsection%
    \noindent #1 \hfill \textit{#2}\par\nopagebreak
    {\small #3}\par\vspace{0.45em}%
}

% --- INÍCIO DO DOCUMENTO ---
\begin{document}

% --- CABEÇALHO ---
\begin{center}
    {\LARGE \textbf{Thiago Buarque}}
    \\[0.15cm]
    {\small Desenvolvedor Backend \textbullet{} .NET/C\# \textbullet{} System Design}
    \\[0.15cm]
    Aracaju, Sergipe
    {\textbullet}
    \href{mailto:thiagobuarque.ma@gmail.com}{thiagobuarque.ma@gmail.com}
    {\textbullet}
    \href{https://www.linkedin.com/in/thiagobqq}{linkedin.com/in/thiagobqq}
    {\textbullet}
    \href{https://github.com/thiagobqq}{github.com/thiagobqq}
    {\textbullet}
    \href{https://thiagobqq.vercel.app/}{thiagobqq.vercel.app}
\end{center}

% --- SEÇÕES ---

\section{Experiência}
    \cventry{Zdoc}{Aracaju, SE (alocado na Desenvolve-se)}{Desenvolvedor Backend .NET}{Set 2025 - Atual}
        \begin{itemize}
            \item Arquiteturei e construí do zero 7 módulos corporativos (RH, Transporte, Transparência, Help Desk e outros) em .NET 9, EF Core e SQL Server, com Clean Architecture em 4 camadas.
            \item Otimizei o tempo de resposta da API de segundos para milissegundos com um pipeline de upload assíncrono para Synology NAS e fila de reprocessamento automático em caso de falha.
            \item Estruturei toda a infraestrutura de deploy do zero (Ubuntu, Docker, Nginx e CI/CD com GitHub Actions), eliminando publicações manuais e automatizando cada atualização.
            \item Integrei login corporativo via LDAP/Identity, autenticação JWT e logs automáticos de auditoria em todo o banco, garantindo rastreabilidade e segurança.
            \item Implementei funcionalidades em tempo real com SignalR e SSE (chat e notificações) e um worker de web scraping para coleta diária de preços de combustível (ANP).
            \item Modelei dados e otimizei consultas em SQL Server/EF Core, sustentando a evolução dos módulos com separação clara de responsabilidades.
        \end{itemize}

% ------

\section{Projetos}
    \projectentry{\href{https://github.com/thiagobqq/Cinema-API}{\textbf{Cinema-API}}}{.NET \textbullet{} Clean Architecture \textbullet{} DDD}{Backend de cinema e venda de ingressos em monolito modular, com módulos isolados de autenticação, filmes, ingressos e pagamentos.}

    \projectentry{\textbf{Vaga Fácil}}{Spring Boot \textbullet{} React \textbullet{} API Pix}{Plataforma freelancer que analisa currículos automaticamente e integra pagamentos via Pix; stack completa em Spring Boot e React (projeto privado).}

    \projectentry{\href{https://github.com/thiagobqq/SearchEngine}{\textbf{SearchEngine}}}{C\# \textbullet{} .NET}{Motor de busca construído em C\# do zero, com indexação eficiente de documentos, estruturas de dados próprias e consulta por relevância.}

    \projectentry{\href{https://github.com/thiagobqq/chat-irl}{\textbf{chat-irl}}}{.NET \textbullet{} SignalR \textbullet{} React \textbullet{} WebSocket}{Chat fullstack em tempo real com SignalR sobre WebSocket: back-end .NET em camadas (Core/Infra/Web) e front-end React com roteamento e páginas dedicadas.}

    \projectentry{\href{https://portal.unit.br/blog/noticias/farmacia-pet-une-saude-animal-tecnologia-e-sustentabilidade-no-conecta-farolandia/}{\textbf{TomPet}}}{React \textbullet{} .NET}{Projeto desenvolvido para a SEMA Aracaju: sistema Farmácia Pet com front-end React e back-end .NET para controle de estoque, prescrições e navegação de operadores.}

% ------

\section{Habilidades}
    \begin{itemize}
        \item \textbf{Back-end:} C\#/.NET 9, ASP.NET Core, Java/Spring Boot, APIs REST
        \item \textbf{Arquitetura:} Clean Architecture, DDD, SOLID, monolito modular, event-driven, mensageria (RabbitMQ), sistemas assíncronos
        \item \textbf{Dados:} SQL Server, EF Core, PostgreSQL, Redis
        \item \textbf{Infra \& DevOps:} Docker, Linux (Ubuntu), Nginx, CI/CD (GitHub Actions), Git
        \item \textbf{IA:} Agentes de IA, programação com IA
        \item \textbf{Front-end (conhecimento):} React, Next.js, TypeScript, Tailwind CSS, Material UI
        \item \textbf{IA:} agentes de IA, programação assistida por IA e Spec-Driven Development
        \item \textbf{Idiomas:} Português (Nativo), Inglês (Avançado/Fluente)
    \end{itemize}

% ------

\section{Educação}
    % Formação mais recente primeiro
    \cventry{Universidade Tiradentes (Unit)}{Aracaju, SE}{Bacharelado em Ciência da Computação}{2023 - 2027 (em andamento)}

\end{document}
